\documentclass[11pt,a4paper]{article}
\usepackage[utf8]{inputenc}
\usepackage[T1]{fontenc}
\usepackage[catalan]{babel}
\usepackage{lmodern}
\usepackage[sfdefault,lining]{FiraSans}
\usepackage[a4paper,top=2.3cm,bottom=2.2cm,left=1.8cm,right=1.8cm,headheight=1.3cm,headsep=0.4cm,footskip=1cm]{geometry}
\usepackage{amsmath,amssymb}
\usepackage{icomma}
\usepackage{graphicx}
\usepackage[table]{xcolor}
\usepackage{tikz}
\usepackage{fontawesome5}
\usepackage{ccicons}
\usepackage{fancyhdr}
\usepackage{lastpage}
\usepackage{tasks}
\usepackage{needspace}
\usepackage[most]{tcolorbox}
\usepackage{array}
\usepackage[hidelinks,pdftitle={Solucions de la fitxa de la unitat 5: Càlcul de derivades. Derivabilitat d'una funció 2n de Batxillerat CCSS},pdfsubject={Matemàtiques Aplicades a les Ciències Socials II},pdfkeywords={derivades, taxa de variació mitjana, definició de derivada, recta tangent, regles de derivació, regla de la cadena, derivades laterals, derivabilitat, funcions a trossos, paràmetres, PAU}]{hyperref}

\definecolor{navy}{HTML}{1B2748}
\definecolor{accent}{HTML}{E8603A}
\definecolor{teal}{HTML}{1F9E90}
\definecolor{soft}{HTML}{EEF1F8}
\definecolor{muted}{HTML}{5D6685}

\setlength{\parindent}{0pt}
\setlength{\parskip}{3pt}
\settasks{label=\alph*), label-format=\bfseries\color{navy}, label-width=1.3em, item-indent=1.8em, column-sep=0.6em, after-item-skip=5pt, before-skip=3pt, after-skip=4pt}

% ---------- capçalera i peu ----------
\newcommand{\llicencia}{\href{https://creativecommons.org/licenses/by-nc-sa/4.0/deed.ca}{\ccbyncsa}\ \ Llicència Creative Commons Reconeixement-NoComercial-CompartirIgual 4.0 Internacional (CC BY-NC-SA 4.0)}
\pagestyle{fancy}
\fancyhf{}
\fancyhead[L]{\raisebox{-0.35\height}{\includegraphics[height=0.95cm]{logo-gva.png}}}
\fancyhead[C]{\small\color{muted}Solucions de la fitxa · Unitat 5: Càlcul de derivades · 2n Batxillerat CCSS}
\fancyhead[R]{\raisebox{-0.35\height}{\includegraphics[height=0.95cm]{logo-vermellar.png}}}
\fancyfoot[L]{\scriptsize\color{muted}\llicencia}
\fancyfoot[R]{\small\color{muted}\thepage}
\renewcommand{\headrulewidth}{0.4pt}
\renewcommand{\headrule}{\color{muted!40}\hrule height \headrulewidth}
\fancypagestyle{primera}{\fancyhf{}\renewcommand{\headrulewidth}{0pt}
  \fancyfoot[L]{\scriptsize\color{muted}\llicencia}
  \fancyfoot[R]{\small\color{muted}\thepage}}

% ---------- nivells i exercicis ----------
\newcommand{\estrelles}[1]{{\color{accent}\ifcase#1\or\faStar\faStar[regular]\faStar[regular]\or\faStar\faStar\faStar[regular]\or\faStar\faStar\faStar\fi}}
\newcounter{exer}
\newcommand{\ex}[1]{\par\needspace{4\baselineskip}\medskip\refstepcounter{exer}\noindent{\color{navy}\bfseries\theexer.}\ \ }

\newcommand{\seccio}[2]{\par\needspace{10\baselineskip}\bigskip
  \begin{tcolorbox}[colback=navy,colframe=navy,coltext=white,arc=2mm,boxrule=0pt,left=3mm,right=3mm,top=1.2mm,bottom=1.2mm]
  {\large\bfseries #1\quad #2}
  \end{tcolorbox}\vspace{-2pt}}


% quadre de recordatori (només a la fitxa d'exercicis)
\newcommand{\recorda}[1]{\par\smallskip\begin{tcolorbox}[colback=accent!7,colframe=accent!55,arc=2mm,boxrule=0.5pt,left=2.5mm,right=2.5mm,top=1mm,bottom=1mm]\small{\color{accent}\faLightbulb[regular]}\ \textbf{Recorda:}\ #1\end{tcolorbox}}
% etiqueta dels exercicis de la PAU d'altres comunitats
\newcommand{\pau}[1]{{\color{teal}\small\bfseries\faUniversity\ PAU · #1.}\ }
\newcommand{\res}[1]{\boxed{#1}}\newcommand{\rest}[1]{\fbox{#1}}

% =====================================================================
% Capçalera FentMates de les fitxes: un color per a cada curs.
% Els colors són els de la pàgina inicial (LEVELS a assets/js/materials.js);
% el logo ve de _eines/logo/pdf/ (python3 _eines/logo/genera.py).
%
% Ús, al final del preàmbul de la unitat:
%   % =====================================================================
% Capçalera FentMates de les fitxes: un color per a cada curs.
% Els colors són els de la pàgina inicial (LEVELS a assets/js/materials.js);
% el logo ve de _eines/logo/pdf/ (python3 _eines/logo/genera.py).
%
% Ús, al final del preàmbul de la unitat:
%   % =====================================================================
% Capçalera FentMates de les fitxes: un color per a cada curs.
% Els colors són els de la pàgina inicial (LEVELS a assets/js/materials.js);
% el logo ve de _eines/logo/pdf/ (python3 _eines/logo/genera.py).
%
% Ús, al final del preàmbul de la unitat:
%   \input{../comu/fentmates.tex}
%   \fmcurs{1eso}
% i, al començament del document:
%   \capcalera{Matemàtiques}{1r ESO}{Fitxa d'exercicis · Unitat 2}{Divisibilitat}
% Opcional abans de \input: \def\fmneed{14} (espai mínim abans de cada secció).
% =====================================================================
\makeatletter
\def\fm@logodir{../../logo/pdf}
\providecommand{\fmneed}{10}
\def\fm@cursnom{}

% clau del curs, color i color del text damunt del color
\newcommand{\fm@curs}[3]{\@namedef{fm@color@#1}{#2}\@namedef{fm@text@#1}{#3}}
\fm@curs{1eso}{E6B84A}{navy}
\fm@curs{2eso}{3B82F6}{white}
\fm@curs{3eso}{FF7A50}{white}
\fm@curs{4eso-a}{A78BFA}{white}
\fm@curs{4eso-b}{6D28D9}{white}
\fm@curs{1batx-mi}{9B1C1C}{white}
\fm@curs{1batx-ccss}{EF6B6B}{white}
\fm@curs{2batx-mii}{1E6B47}{white}
\fm@curs{2batx-ccss}{2BB3A3}{white}

\newcommand{\fmcurs}[1]{%
  \@ifundefined{fm@color@#1}{\PackageError{fentmates}{Curs desconegut: #1}{}}{}%
  \edef\fm@tmp{\noexpand\definecolor{curs}{HTML}{\@nameuse{fm@color@#1}}%
               \noexpand\colorlet{curstext}{\@nameuse{fm@text@#1}}}\fm@tmp
  \colorlet{cursfosc}{curs!72!black}%
  \def\fm@logo{\fm@logodir/vertical-#1.pdf}%
  \def\fm@marca{\fm@logodir/marca-#1.pdf}}

% ---------- capçalera de la primera pàgina ----------
% #1 matèria (petit), #2 curs, #3 tipus de document i unitat, #4 títol
\newcommand{\capcalera}[4]{%
  \gdef\fm@cursnom{#2}%
  \par\noindent
  \begin{tcolorbox}[colback=curs,colframe=curs,boxrule=0pt,arc=2.2mm,width=3.5cm,height=2.15cm,
      valign=center,left=3mm,right=2mm,top=1mm,bottom=1mm,nobeforeafter,box align=center]
    \color{curstext}{\footnotesize #1}\par\vspace{1pt}{\fontsize{19}{21}\selectfont\bfseries #2\par}%
  \end{tcolorbox}\hspace{5mm}%
  \begin{minipage}[c]{\dimexpr\linewidth-3.5cm-5mm-2.9cm\relax}
    {\color{cursfosc}\small\bfseries\MakeUppercase{#3}}\par\vspace{2pt}
    {\color{navy}\fontsize{20}{23}\selectfont\bfseries #4\par}
  \end{minipage}\hfill
  \begin{minipage}[c]{2.6cm}\raggedleft\includegraphics[width=2.6cm]{\fm@logo}\end{minipage}\par
  \vspace{2pt}}

% ---------- barres de secció del color del curs ----------
\renewcommand{\seccio}[2]{\par\needspace{\fmneed\baselineskip}\bigskip
  \begin{tcolorbox}[colback=curs,colframe=curs,coltext=curstext,arc=2mm,boxrule=0pt,left=3mm,right=3mm,top=1.2mm,bottom=1.2mm]
  {\large\bfseries #1\quad #2}
  \end{tcolorbox}\vspace{-2pt}}

% ---------- capçalera i peu de les pàgines ----------
% primera pàgina: logos institucionals al peu
\fancypagestyle{primera}{\fancyhf{}\renewcommand{\headrulewidth}{0pt}
  \fancyfoot[L]{\raisebox{-0.32\height}{\includegraphics[height=0.78cm]{logo-gva.png}}\hspace{5mm}%
    \raisebox{-0.32\height}{\includegraphics[height=0.78cm]{logo-vermellar.png}}}
  \fancyfoot[R]{\scriptsize\color{muted}\href{https://creativecommons.org/licenses/by-nc-sa/4.0/deed.ca}{\ccbyncsa}\ \ CC BY-NC-SA 4.0 · fentmates.org\qquad\small\thepage}}
% la resta: el símbol a l'esquerra i el curs a la dreta, amb el títol al centre com abans
\fancyhead[L]{\raisebox{-0.3\height}{\includegraphics[height=0.62cm]{\fm@marca}}}
\fancyhead[R]{\small\bfseries\color{cursfosc}\fm@cursnom}
\renewcommand{\headrule}{\color{curs!45}\hrule height \headrulewidth}
\makeatother

%   \fmcurs{1eso}
% i, al començament del document:
%   \capcalera{Matemàtiques}{1r ESO}{Fitxa d'exercicis · Unitat 2}{Divisibilitat}
% Opcional abans de \input: \def\fmneed{14} (espai mínim abans de cada secció).
% =====================================================================
\makeatletter
\def\fm@logodir{../../logo/pdf}
\providecommand{\fmneed}{10}
\def\fm@cursnom{}

% clau del curs, color i color del text damunt del color
\newcommand{\fm@curs}[3]{\@namedef{fm@color@#1}{#2}\@namedef{fm@text@#1}{#3}}
\fm@curs{1eso}{E6B84A}{navy}
\fm@curs{2eso}{3B82F6}{white}
\fm@curs{3eso}{FF7A50}{white}
\fm@curs{4eso-a}{A78BFA}{white}
\fm@curs{4eso-b}{6D28D9}{white}
\fm@curs{1batx-mi}{9B1C1C}{white}
\fm@curs{1batx-ccss}{EF6B6B}{white}
\fm@curs{2batx-mii}{1E6B47}{white}
\fm@curs{2batx-ccss}{2BB3A3}{white}

\newcommand{\fmcurs}[1]{%
  \@ifundefined{fm@color@#1}{\PackageError{fentmates}{Curs desconegut: #1}{}}{}%
  \edef\fm@tmp{\noexpand\definecolor{curs}{HTML}{\@nameuse{fm@color@#1}}%
               \noexpand\colorlet{curstext}{\@nameuse{fm@text@#1}}}\fm@tmp
  \colorlet{cursfosc}{curs!72!black}%
  \def\fm@logo{\fm@logodir/vertical-#1.pdf}%
  \def\fm@marca{\fm@logodir/marca-#1.pdf}}

% ---------- capçalera de la primera pàgina ----------
% #1 matèria (petit), #2 curs, #3 tipus de document i unitat, #4 títol
\newcommand{\capcalera}[4]{%
  \gdef\fm@cursnom{#2}%
  \par\noindent
  \begin{tcolorbox}[colback=curs,colframe=curs,boxrule=0pt,arc=2.2mm,width=3.5cm,height=2.15cm,
      valign=center,left=3mm,right=2mm,top=1mm,bottom=1mm,nobeforeafter,box align=center]
    \color{curstext}{\footnotesize #1}\par\vspace{1pt}{\fontsize{19}{21}\selectfont\bfseries #2\par}%
  \end{tcolorbox}\hspace{5mm}%
  \begin{minipage}[c]{\dimexpr\linewidth-3.5cm-5mm-2.9cm\relax}
    {\color{cursfosc}\small\bfseries\MakeUppercase{#3}}\par\vspace{2pt}
    {\color{navy}\fontsize{20}{23}\selectfont\bfseries #4\par}
  \end{minipage}\hfill
  \begin{minipage}[c]{2.6cm}\raggedleft\includegraphics[width=2.6cm]{\fm@logo}\end{minipage}\par
  \vspace{2pt}}

% ---------- barres de secció del color del curs ----------
\renewcommand{\seccio}[2]{\par\needspace{\fmneed\baselineskip}\bigskip
  \begin{tcolorbox}[colback=curs,colframe=curs,coltext=curstext,arc=2mm,boxrule=0pt,left=3mm,right=3mm,top=1.2mm,bottom=1.2mm]
  {\large\bfseries #1\quad #2}
  \end{tcolorbox}\vspace{-2pt}}

% ---------- capçalera i peu de les pàgines ----------
% primera pàgina: logos institucionals al peu
\fancypagestyle{primera}{\fancyhf{}\renewcommand{\headrulewidth}{0pt}
  \fancyfoot[L]{\raisebox{-0.32\height}{\includegraphics[height=0.78cm]{logo-gva.png}}\hspace{5mm}%
    \raisebox{-0.32\height}{\includegraphics[height=0.78cm]{logo-vermellar.png}}}
  \fancyfoot[R]{\scriptsize\color{muted}\href{https://creativecommons.org/licenses/by-nc-sa/4.0/deed.ca}{\ccbyncsa}\ \ CC BY-NC-SA 4.0 · fentmates.org\qquad\small\thepage}}
% la resta: el símbol a l'esquerra i el curs a la dreta, amb el títol al centre com abans
\fancyhead[L]{\raisebox{-0.3\height}{\includegraphics[height=0.62cm]{\fm@marca}}}
\fancyhead[R]{\small\bfseries\color{cursfosc}\fm@cursnom}
\renewcommand{\headrule}{\color{curs!45}\hrule height \headrulewidth}
\makeatother

%   \fmcurs{1eso}
% i, al començament del document:
%   \capcalera{Matemàtiques}{1r ESO}{Fitxa d'exercicis · Unitat 2}{Divisibilitat}
% Opcional abans de \input: \def\fmneed{14} (espai mínim abans de cada secció).
% =====================================================================
\makeatletter
\def\fm@logodir{../../logo/pdf}
\providecommand{\fmneed}{10}
\def\fm@cursnom{}

% clau del curs, color i color del text damunt del color
\newcommand{\fm@curs}[3]{\@namedef{fm@color@#1}{#2}\@namedef{fm@text@#1}{#3}}
\fm@curs{1eso}{E6B84A}{navy}
\fm@curs{2eso}{3B82F6}{white}
\fm@curs{3eso}{FF7A50}{white}
\fm@curs{4eso-a}{A78BFA}{white}
\fm@curs{4eso-b}{6D28D9}{white}
\fm@curs{1batx-mi}{9B1C1C}{white}
\fm@curs{1batx-ccss}{EF6B6B}{white}
\fm@curs{2batx-mii}{1E6B47}{white}
\fm@curs{2batx-ccss}{2BB3A3}{white}

\newcommand{\fmcurs}[1]{%
  \@ifundefined{fm@color@#1}{\PackageError{fentmates}{Curs desconegut: #1}{}}{}%
  \edef\fm@tmp{\noexpand\definecolor{curs}{HTML}{\@nameuse{fm@color@#1}}%
               \noexpand\colorlet{curstext}{\@nameuse{fm@text@#1}}}\fm@tmp
  \colorlet{cursfosc}{curs!72!black}%
  \def\fm@logo{\fm@logodir/vertical-#1.pdf}%
  \def\fm@marca{\fm@logodir/marca-#1.pdf}}

% ---------- capçalera de la primera pàgina ----------
% #1 matèria (petit), #2 curs, #3 tipus de document i unitat, #4 títol
\newcommand{\capcalera}[4]{%
  \gdef\fm@cursnom{#2}%
  \par\noindent
  \begin{tcolorbox}[colback=curs,colframe=curs,boxrule=0pt,arc=2.2mm,width=3.5cm,height=2.15cm,
      valign=center,left=3mm,right=2mm,top=1mm,bottom=1mm,nobeforeafter,box align=center]
    \color{curstext}{\footnotesize #1}\par\vspace{1pt}{\fontsize{19}{21}\selectfont\bfseries #2\par}%
  \end{tcolorbox}\hspace{5mm}%
  \begin{minipage}[c]{\dimexpr\linewidth-3.5cm-5mm-2.9cm\relax}
    {\color{cursfosc}\small\bfseries\MakeUppercase{#3}}\par\vspace{2pt}
    {\color{navy}\fontsize{20}{23}\selectfont\bfseries #4\par}
  \end{minipage}\hfill
  \begin{minipage}[c]{2.6cm}\raggedleft\includegraphics[width=2.6cm]{\fm@logo}\end{minipage}\par
  \vspace{2pt}}

% ---------- barres de secció del color del curs ----------
\renewcommand{\seccio}[2]{\par\needspace{\fmneed\baselineskip}\bigskip
  \begin{tcolorbox}[colback=curs,colframe=curs,coltext=curstext,arc=2mm,boxrule=0pt,left=3mm,right=3mm,top=1.2mm,bottom=1.2mm]
  {\large\bfseries #1\quad #2}
  \end{tcolorbox}\vspace{-2pt}}

% ---------- capçalera i peu de les pàgines ----------
% primera pàgina: logos institucionals al peu
\fancypagestyle{primera}{\fancyhf{}\renewcommand{\headrulewidth}{0pt}
  \fancyfoot[L]{\raisebox{-0.32\height}{\includegraphics[height=0.78cm]{logo-gva.png}}\hspace{5mm}%
    \raisebox{-0.32\height}{\includegraphics[height=0.78cm]{logo-vermellar.png}}}
  \fancyfoot[R]{\scriptsize\color{muted}\href{https://creativecommons.org/licenses/by-nc-sa/4.0/deed.ca}{\ccbyncsa}\ \ CC BY-NC-SA 4.0 · fentmates.org\qquad\small\thepage}}
% la resta: el símbol a l'esquerra i el curs a la dreta, amb el títol al centre com abans
\fancyhead[L]{\raisebox{-0.3\height}{\includegraphics[height=0.62cm]{\fm@marca}}}
\fancyhead[R]{\small\bfseries\color{cursfosc}\fm@cursnom}
\renewcommand{\headrule}{\color{curs!45}\hrule height \headrulewidth}
\makeatother

\fmcurs{2batx-ccss}

\begin{document}
\thispagestyle{primera}

\capcalera{Matemàtiques CCSS}{2n Batx}{Solucions de la fitxa · Unitat 5}{Càlcul de derivades. Derivabilitat d'una funció}
\vspace{8pt}

\begin{tcolorbox}[colback=soft,colframe=soft,arc=2mm,left=3mm,right=3mm,top=1.5mm,bottom=1.5mm]
\small Intenta resoldre cada exercici abans de mirar-ne la solució. Si no t'ix el mateix resultat, repassa els passos o comprova'l amb la calculadora de derivades dels apunts interactius de la unitat. Una derivada pot estar escrita de maneres diferents i ser igualment correcta.
\end{tcolorbox}

% =====================================================
\seccio{1}{Taxa de variació mitjana i interpretació de la derivada}

\ex{1} $f(0)=1,\ f(1)=-2,\ f(2)=-3,\ f(4)=1,\ f(5)=6$.\begin{tasks}(3)
  \task $\frac{-3-1}{2}=\res{-2}$
  \task $\frac{1-(-2)}{3}=\res{1}$
  \task $\frac{6-(-3)}{3}=\res{3}$
\end{tasks}

\ex{1} $\text{TVM}[1,4]=\frac{2-1}{3}=\res{\frac13}$; \ $\text{TVM}[4,9]=\frac{3-2}{5}=\res{\frac15}$. Creix més ràpidament en \rest{$[1,4]$}.

\ex{2} \label{ex:cost}$C(100)=800$, $C(200)=1600$, $C(300)=2600$. \ $\text{TVM}[100,200]=\frac{800}{100}=\res{8}$ €/unitat; \ $\text{TVM}[200,300]=\frac{1000}{100}=\res{10}$ €/unitat: és el cost mitjà de cada unitat de més en cada tram, i va augmentant.

\ex{1} \begin{tasks}(1)
  \task $y-2=-4(x-3)\Rightarrow\res{y=-4x+14}$
  \task $\res{g'(1)=-2}$ (pendent) i $\res{g(1)=-2+7=5}$.
  \task El pendent és 0: la tangent és \rest{horitzontal}, $y=f(a)$.
\end{tasks}


% =====================================================
\seccio{2}{Derivada d'una funció per definició}

\ex{2} \begin{tasks}(1)
  \task $\frac{3(1+h)^2-(1+h)-2}{h}=\frac{5h+3h^2}{h}=5+3h\to\res{5}$
  \task $f(-1+h)=(-1+h)^3-2(-1+h)=1+h-3h^2+h^3$, \ $f(-1)=1$: \ $\frac{h-3h^2+h^3}{h}=1-3h+h^2\to\res{1}$
  \task $\frac{\frac{2}{2+h}-1}{h}=\frac{\frac{-h}{2+h}}{h}=\frac{-1}{2+h}\to\res{-\frac12}$
  \task $\frac{\sqrt{9+2h}-3}{h}=\frac{9+2h-9}{h(\sqrt{9+2h}+3)}=\frac{2}{\sqrt{9+2h}+3}\to\res{\frac13}$
\end{tasks}

\ex{2} \pau{Navarra, EvAU extraordinària 2023}$g(1+h)=-(1+h)^2+4(1+h)+3=6+2h-h^2$, \ $g(1)=6$. \ $g'(1)=\displaystyle\lim_{h\to0}\frac{2h-h^2}{h}=\lim_{h\to0}(2-h)=\res{2}$.

\ex{2} \begin{tasks}(1)
  \task $\frac{(x+h)^2-3(x+h)+2-(x^2-3x+2)}{h}=\frac{2xh+h^2-3h}{h}=2x+h-3\to\res{2x-3}$
  \task $\frac{\frac{1}{x+h+2}-\frac{1}{x+2}}{h}=\frac{-1}{(x+h+2)(x+2)}\to\res{-\frac{1}{(x+2)^2}}$
  \task $\frac{\sqrt{x+h}-\sqrt x}{h}=\frac{1}{\sqrt{x+h}+\sqrt x}\to\res{\frac{1}{2\sqrt x}}$
\end{tasks}

\ex{2} $\frac{(3+h)^2-9}{h}=6+h\to f'(3)=6$. \ $f(3)=9$: \ $y-9=6(x-3)\Rightarrow\res{y=6x-9}$.


% =====================================================
\seccio{3}{Propietats de la derivada: suma, producte i quocient}

\ex{1} \begin{tasks}(2)
  \task $\res{20x^3-9x^2+2}$
  \task $\res{x^4-2x^2+1}$
  \task $\frac{6x-2}{4}=\res{\frac{3x-1}{2}}$
  \task $\res{-6x^5+\frac12}$
\end{tasks}

\ex{1} \begin{tasks}(1)
  \task $2(x^2+3)+(2x-1)\cdot2x=\res{6x^2-2x+6}$
  \task $(3x^2+1)(1-x^2)+(x^3+x)(-2x)=\res{1-5x^4}$
  \task $f(x)=x^4-x^2\Rightarrow\res{4x^3-2x}$
\end{tasks}

\ex{2} \begin{tasks}(2)
  \task $\frac{(x-2)-(x+3)}{(x-2)^2}=\res{\frac{-5}{(x-2)^2}}$
  \task $\frac{4x(x+1)-2x^2}{(x+1)^2}=\res{\frac{2x^2+4x}{(x+1)^2}}$
  \task $\frac{2x(x^2+1)-(x^2-1)2x}{(x^2+1)^2}=\res{\frac{4x}{(x^2+1)^2}}$
  \task $\res{\frac{-6x}{(x^2-4)^2}}$
\end{tasks}

\ex{1} $\res{f'(x)=4x^3-6x^2+5}$, \ $\res{f''(x)=12x^2-12x}$, \ $\res{f''(1)=0}$.

\ex{2} $\res{C'(x)=0{,}02x+5}$, \ $\res{C'(100)=7}$: quan ja se'n fabriquen 100, fabricar-ne una més costa aproximadament 7 €. La TVM en $[100,200]$ (8 €) era una mitjana de tot el tram; la derivada mesura el ritme \emph{just} en $x=100$.


% =====================================================
\seccio{4}{Derivades de les funcions elementals}

\ex{1} \begin{tasks}(3)
  \task $x^{3/2}\to\res{\frac32\sqrt x}$
  \task $x^{-2}\to\res{\frac{-2}{x^3}}$
  \task $x^{1/3}\to\res{\frac{1}{3\sqrt[3]{x^2}}}$
  \task $6x^{-1/2}\to\res{\frac{-3}{x\sqrt x}}$
  \task $x^{5/2}\to\res{\frac52x\sqrt x}$
  \task $\frac23x^{-3}\to\res{\frac{-2}{x^4}}$
\end{tasks}

\ex{1} \begin{tasks}(3)
  \task $\res{2e^x-\frac3x}$
  \task $\res{4^x\ln4}$
  \task $\res{\frac{1}{x\ln3}}$
  \task $e^x+xe^x=\res{(x+1)e^x}$
  \task $\frac{e^xx-e^x}{x^2}=\res{\frac{e^x(x-1)}{x^2}}$
  \task $\res{2x\ln x+x}$
\end{tasks}

\ex{2} \begin{tasks}(2)
  \task $\res{3\cos x+2\sin x}$
  \task $\res{\sin x+x\cos x}$
  \task $1+\tan^2x-1=\res{\tan^2x}$
  \task $\res{\cos^2x-\sin^2x}$
  \task $\frac{1}{1+x^2}-1=\res{\frac{-x^2}{1+x^2}}$
  \task $\res{\arcsin x+\frac{x}{\sqrt{1-x^2}}}$
\end{tasks}


% =====================================================
\seccio{5}{Regla de la cadena}

\ex{2} \begin{tasks}(2)
  \task $4(3x^2-1)^3\cdot6x=\res{24x(3x^2-1)^3}$
  \task $\frac{2x-2}{2\sqrt{x^2-2x}}=\res{\frac{x-1}{\sqrt{x^2-2x}}}$
  \task $-3(x^2+1)^{-4}\cdot2x=\res{\frac{-6x}{(x^2+1)^4}}$
  \task $\frac13(2x+5)^{-2/3}\cdot2=\res{\frac{2}{3\sqrt[3]{(2x+5)^2}}}$
\end{tasks}

\ex{2} \begin{tasks}(3)
  \task $\res{3e^{3x-1}}$
  \task $\res{2x\cdot2^{x^2}\ln2}$
  \task $\res{\frac{2x-3}{x^2-3x}}$
  \task $\frac{5}{5x\ln10}=\res{\frac{1}{x\ln10}}$
  \task $\frac12\ln x\to\res{\frac{1}{2x}}$
  \task $\res{\frac{e^{\sqrt x}}{2\sqrt x}}$
\end{tasks}

\ex{2} \begin{tasks}(3)
  \task $\res{5\cos(5x+1)}$
  \task $\res{-2x\sin(x^2)}$
  \task $\res{\frac{3}{\cos^2(3x)}}$
  \task $\res{3\sin^2x\cos x}$
  \task $\res{\frac{2}{1+4x^2}}$
  \task $\res{\frac{2x}{\sqrt{1-x^4}}}$
\end{tasks}

\ex{2} \begin{tasks}(1)
  \task $e^{-x^2}+x\,e^{-x^2}(-2x)=\res{(1-2x^2)\,e^{-x^2}}$
  \task $\ln(x+1)-\ln(x^2+1)\Rightarrow\res{\frac{1}{x+1}-\frac{2x}{x^2+1}}=\frac{1-2x-x^2}{(x+1)(x^2+1)}$
  \task $2(x-1)(x+2)^3+3(x-1)^2(x+2)^2=\res{(x-1)(x+2)^2(5x+1)}$
  \task $\frac{e^x(e^x+1)-e^xe^x}{(e^x+1)^2}=\res{\frac{e^x}{(e^x+1)^2}}$
\end{tasks}

\ex{2} \pau{Múrcia, EBAU setembre 2020}\begin{tasks}(2)
  \task $\res{2x\ln x+\frac{x^2-2}{x}}$
  \task $\res{12x^2e^{4x^3+2}}$
\end{tasks}

\ex{2} \pau{Múrcia, EBAU juny 2018}\begin{tasks}(2)
  \task $2x\ln(x^2)+x^2\cdot\frac{2x}{x^2}=\res{2x\ln(x^2)+2x}$
  \task $\res{(2x-3)\,e^{-3x+x^2}}$
\end{tasks}

\ex{2} \pau{Andalusia, PEvAU ordinària 2021}\begin{tasks}(2)
  \task $\frac{1}{x-1}-\frac{1}{x+1}=\res{\frac{2}{x^2-1}}$
  \task $3x^2e^{2x^2}+x^3e^{2x^2}\cdot4x=\res{e^{2x^2}(3x^2+4x^4)}$
\end{tasks}

\ex{2} \pau{Andalusia, ordinària (suplent) 2025}\begin{tasks}(1)
  \task $g'(x)=\frac{2x}{x^2-1}-\frac{2x}{x^2+1}=\res{\frac{4x}{(x^2-1)(x^2+1)}}$
  \task $h'(x)=2e^{x^2-x}+(2x-1)^2e^{x^2-x}=\res{(4x^2-4x+3)\,e^{x^2-x}}$
\end{tasks}

\ex{3} \pau{Andalusia, PEvAU ordinària 2024}\begin{tasks}(1)
  \task $f'(x)=6x(x^2+2)^2e^{-2x}-2(x^2+2)^3e^{-2x}=\res{-2(x^2-3x+2)(x^2+2)^2e^{-2x}}$
  \task $g'(x)=\dfrac{\frac{-3x^2}{1-x^3}(1-2x^2)^2-\ln(1-x^3)\cdot2(1-2x^2)(-4x)}{(1-2x^2)^4}=\res{\dfrac{-3x^2(1-2x^2)+8x(1-x^3)\ln(1-x^3)}{(1-2x^2)^3(1-x^3)}}$
\end{tasks}

\ex{2} \pau{Múrcia, EBAU setembre 2018}\begin{tasks}(1)
  \task $f'(x)=15x^2e^{2x}+10x^3e^{2x}-\frac{2x}{(x^2+1)^2}\Rightarrow\res{f'(0)=0}$
  \task $y-1=0\cdot(x-0)\Rightarrow\res{y=1}$
\end{tasks}

\ex{2} \begin{tasks}(1)
  \task $f(1)=\frac12$, $f'(x)=\frac{-1}{(x+1)^2}$, $f'(1)=-\frac14$: \ $y-\frac12=-\frac14(x-1)\Rightarrow\res{y=-0{,}25x+0{,}75}$
  \task $f(-2)=3$, $f'(x)=\frac{x}{\sqrt{5+x^2}}$, $f'(-2)=-\frac23$: \ $y-3=-\frac23(x+2)\Rightarrow\res{y=-\frac23x+\frac53}$
  \task $f(1)=-\frac13$, $f'(x)=\frac{-x^2-4}{(x^2-4)^2}$, $f'(1)=-\frac59$: \ $y+\frac13=-\frac59(x-1)\Rightarrow\res{5x+9y-2=0}$
\end{tasks}

\ex{3} \pau{Balears, PBAU extraordinària 2021}\begin{tasks}(1)
  \task La recta té pendent $-\frac34$. \ $f'(x)=-\frac{3}{x^2}=-\frac34\Rightarrow x^2=4\Rightarrow x=\pm2$: \ punts $\res{(-2;\ 6{,}5)}$ i $\res{(2;\ 9{,}5)}$.
  \task En $x=-2$: $\res{y=-\frac34x+5}$; \ en $x=2$: $\res{y=-\frac34x+11}$.
\end{tasks}

\ex{2} \begin{tasks}(1)
  \task $f'(x)=a-\frac{b}{x^2}$: \ $2a+\frac b2=4$, \ $a-\frac b4=0\Rightarrow\res{a=1,\ b=4}$
  \task $f(0)=b=1$; \ $f'(x)=\frac{a(x^2+1)-(ax+b)2x}{(x^2+1)^2}\Rightarrow f'(0)=a=1$: \ $\res{a=b=1}$
  \task $f(0)=b=2$; \ $f'(x)=3x^2\ln(2x+5)+\frac{2x^3}{2x+5}+a\Rightarrow f'(0)=a=1$: \ $\res{a=1,\ b=2}$
  \task $h'(1)=3+2a+3=0\Rightarrow a=-3$; \ $h(1)=1-3+3+b=2\Rightarrow b=1$: \ $\res{a=-3,\ b=1}$
\end{tasks}


% =====================================================
\seccio{6}{Continuïtat i derivabilitat. Derivades laterals}

\ex{2} \begin{tasks}(1)
  \task $f'(-1^-)=-1$, $f'(-1^+)=1$: \rest{no derivable} (punt anguló).
  \task Prop de 3: $9-x^2$ per l'esquerra i $x^2-9$ per la dreta: $f'(3^-)=-6$, $f'(3^+)=6$: \rest{no derivable}.
  \task Prop de 1: $x(1-x)$ per l'esquerra i $x(x-1)$ per la dreta: $f'(1^-)=1-2=-1$, $f'(1^+)=2-1=1$: \rest{no derivable}.
  \task $\displaystyle\lim_{h\to0}\frac{\sqrt[3]{h^2}}{h}=\lim_{h\to0}\frac{1}{\sqrt[3]{h}}$: $-\infty$ per l'esquerra i $+\infty$ per la dreta: \rest{no derivable}.
\end{tasks}

\ex{1} \begin{tasks}(1)
  \task \rest{Falsa}: $|x|$ és contínua en 0 però no hi és derivable.
  \task \rest{Vertadera}: si fóra derivable, seria contínua.
  \task \rest{Falsa}: $|x|$ no és derivable en 0 i sí que és contínua.
  \task \rest{Falsa}: cal que, a més, siga contínua. Per exemple, $f(x)=x$ si $x<0$ i $f(x)=x+1$ si $x\geq0$: els dos trossos tenen derivada 1, però hi ha un salt en $x=0$.
\end{tasks}

\ex{2} En $x=2$ hi ha un salt: \rest{no és contínua} i, per tant, \rest{no és derivable}. En $x=-2$ és contínua, però hi ha un punt anguló (pendents $1$ i $-1$): \rest{no és derivable}. En $x=0$ és contínua i derivable (unió suau). $f$ és contínua en $\res{\mathbb R-\{2\}}$ i derivable en $\res{\mathbb R-\{-2,2\}}$.


% =====================================================
\seccio{7}{Derivabilitat de funcions a trossos}

\ex{2} \begin{tasks}(1)
  \task $f$: \ $\lim_{x\to0^-}e^x=1=f(0)$: contínua. \ $f'(x)=e^x$ si $x<0$ i $1$ si $x>0$: $f'(0^-)=1=f'(0^+)$. \ \rest{Contínua i derivable en $\mathbb R$}.
  \task $g$: \ $\lim_{x\to1^-}(3-x^2)=2=g(1)$: contínua ($\frac2x$ només falla en $0$, fora de $x\geq1$). \ $g'(1^-)=-2$, \ $g'(1^+)=-\frac{2}{1^2}=-2$. \ \rest{Contínua i derivable en $\mathbb R$}.
\end{tasks}

\ex{2} \pau{Andalusia, PEvAU extraordinària 2024}Trossos polinòmics. $x=4$: $f(4)=3$, $\lim_{x\to4^-}f=-16+16+3=3$: contínua en $\mathbb R$. \ $f'(x)=-2x+4$ si $x<4$, $2$ si $x>4$: $f'(4^-)=-4\neq2=f'(4^+)$. \ \rest{Contínua en $\mathbb R$ i derivable en $\mathbb R-\{4\}$}.

\ex{2} \pau{Andalusia, ordinària (suplent) 2025}$f(-2)=10-5=5$, \ $\lim_{x\to-2^+}(x^2+1)=5$: \rest{contínua} en $x=-2$. \ $f'(-2^-)=\frac52$, \ $f'(-2^+)=2\cdot(-2)=-4$: \rest{no és derivable} en $x=-2$.

\ex{2} \pau{Andalusia, PEvAU ordinària 2024}$t=10$: $v(10)=80$, $\lim_{t\to10^+}v=-100+320-140=80$: contínua. \ $v'(t)=2t-8$ si $0<t<10$ i $-2t+32$ si $10<t<24$: $v'(10^-)=12=v'(10^+)$. \ \rest{Contínua en $[0,24]$ i derivable en $(0,24)$}.

\ex{2} \pau{Navarra, EvAU ordinària 2024}\begin{tasks}(1)
  \task $x=-1$: $\lim_{x\to-1^-}f=3$, $\lim_{x\to-1^+}f=5=f(-1)$: \rest{salt finit}. \ $x=1$: $3=3=f(1)$: contínua. \ Contínua en $\res{\mathbb R-\{-1\}}$.
  \task En $x=-1$ no és derivable (no és contínua). \ $f'(1^-)=-1$, $f'(1^+)=2-6=-4$: no derivable. \ Derivable en $\res{\mathbb R-\{-1,1\}}$.
\end{tasks}

\ex{2} \pau{Andalusia, PEvAU extraordinària 2021}$\lim_{x\to0^-}2^{x+1}=2$, \ $\lim_{x\to0^+}(x^2-2x)=0$: no és contínua en $x=0$ i, per tant, \rest{no és derivable en $x=0$}. Els trossos són contínus i derivables: $f$ és contínua i derivable en $\res{\mathbb R-\{0\}}$.

\ex{3} \pau{Andalusia, PEvAU extraordinària 2021}$t=1{,}8$: $f(1{,}8)=0{,}06$ i $\lim_{t\to1{,}8^+}f=0{,}06$; \ $t=5$: $f(5)=0{,}38$ i $\lim_{t\to5^+}f=0{,}38$: \rest{contínua en $[0{,}2;\ 10]$}. \ $f'(1{,}8^-)=-1{,}6\neq0{,}1=f'(1{,}8^+)$; \ $f'(5^-)=0{,}1\neq3{,}3=f'(5^+)$. \ Derivable en $\res{(0{,}2;\ 1{,}8)\cup(1{,}8;\ 5)\cup(5;\ 10)}$.

\ex{3} \begin{tasks}(1)
  \task $f(x)=-2x$ si $x<-1$; \ $2$ si $-1\leq x\leq1$; \ $2x$ si $x>1$. Contínua; $f'(-1^-)=-2\neq0=f'(-1^+)$ i $f'(1^-)=0\neq2=f'(1^+)$: derivable en $\res{\mathbb R-\{-1,1\}}$.
  \task $g(x)=x^2-4x$ si $x<0$ o $x>4$; \ $-x^2+4x$ si $0\leq x\leq4$. Contínua; $g'(0^-)=-4\neq4=g'(0^+)$ i $g'(4^-)=-4\neq4=g'(4^+)$: derivable en $\res{\mathbb R-\{0,4\}}$.
\end{tasks}


% =====================================================
\seccio{8}{Funcions a trossos amb paràmetres}

\ex{2} \pau{Andalusia, PEvAU ordinària 2022}Continuïtat en $x=1$: $3=a+b+2$. \ Derivabilitat: $f'(1^-)=6$, $f'(1^+)=2a+b\Rightarrow2a+b=6$. \ $\begin{cases}a+b=1\\2a+b=6\end{cases}\Rightarrow\res{a=5,\ b=-4}$.

\ex{2} \pau{Andalusia, PEvAU extraordinària 2020}Continuïtat en $x=0$: $2-a=a+b$. \ $f'(x)=\frac{-a}{(x-1)^2}$ si $x<0$ i $be^x$ si $x>0$: $-a=b$. \ Substituint: $2-a=0\Rightarrow\res{a=2,\ b=-2}$ ($x=1$ no és de $x<0$).

\ex{2} \pau{Balears, PBAU ordinària 2022}Continuïtat: $2=2$, contínua per a qualsevol $a$. \ $f'(0^-)=-3$, $f'(0^+)=a\Rightarrow\res{a=-3}$.

\ex{2} \pau{Balears, PBAU ordinària 2021}Continuïtat: $2=e^0+1=2$, per a qualsevol $a$. \ $f'(0^-)=-3$, $f'(0^+)=a\,e^{0}=a\Rightarrow\res{a=-3}$.

\ex{2} \pau{Balears, PBAU extraordinària 2021}\begin{tasks}(1)
  \task $x=-2$: $8+a=4-5\Rightarrow\res{a=-9}$. \ ($x=1$: $-4=-4$ sempre.)
  \task Per a $a=1$ no és contínua en $x=-2$ ($9\neq-1$): no és derivable. En $x=1$: $f'(1^-)=2\neq-7=f'(1^+)$. \ \rest{No és derivable en $x=-2$ ni en $x=1$}.
\end{tasks}

\ex{3} \pau{Aragó, EvAU ordinària 2022}\begin{tasks}(1)
  \task $x=0$: $2\sqrt3=\sqrt b\Rightarrow b=12$. \ $x=3$: $\sqrt{3a+12}=\frac72-\frac12=3\Rightarrow\res{a=-1,\ b=12}$.
  \task $f'(x)=\frac{-1}{2\sqrt{12-x}}$ si $0<x<3$. \ $f'(0^-)=0\neq-\frac{1}{2\sqrt{12}}=f'(0^+)$: \rest{no derivable en $x=0$}. \ $f'(3^-)=-\frac16=f'(3^+)$: \rest{derivable en $x=3$}.
\end{tasks}

\ex{2} \pau{Andalusia, ordinària (suplent) 2025}$f(0)=\res{10}$. \ Continuïtat: $-6{,}25+5+10=8{,}75=6{,}25+2{,}5a+b$. \ Derivabilitat: $-2\cdot2{,}5+2=-3=5+a\Rightarrow a=-8$, \ i $b=8{,}75-6{,}25+20=22{,}5$: \ $\res{a=-8,\ b=22{,}5}$.

\ex{2} \begin{tasks}(1)
  \task $f$: \ $4a+2b+1=3$ i $4a+b=3\Rightarrow\res{a=1,\ b=-1}$.
  \task $g$: \ $1=b$ i $g'(0^-)=2e^0=2=a\Rightarrow\res{a=2,\ b=1}$.
\end{tasks}


\end{document}
